
\subsection{RQ5: Complexity-Guided Routing System Evaluation}
\label{sec:complexity-routing}

To demonstrate the practical effectiveness of our RL+LLM approach on computationally challenging constraints, we developed a complexity-guided routing system that identifies difficult constraints using static analysis features, without requiring pre-solving.

\textbf{Experimental Design}

Our improved experimental design addresses the key limitation of prediction-based routing by using a \textit{constraint complexity analyzer} that operates on syntactic and structural features:

\begin{enumerate}
    \item \textbf{Static Complexity Analysis}: We extract features including:
    \begin{itemize}
        \item \textit{Nonlinear operations}: Modular arithmetic (\%), power operations (\^), variable interactions
        \item \textit{Structural complexity}: Nested parentheses, operator chains, constraint length
        \item \textit{Variable characteristics}: Variable count, interaction patterns, degree of polynomial terms
    \end{itemize}
    
    \item \textbf{Complexity-Based Classification}: Constraints are classified into three categories:
    \begin{itemize}
        \item \textit{Easy} (complexity score < 5.0): Simple linear constraints with few variables
        \item \textit{Medium} (5.0 ≤ score < 15.0): Moderate nonlinearity or variable interactions
        \item \textit{Hard} (score ≥ 15.0): High nonlinearity, modular arithmetic, or complex interactions
    \end{itemize}
    
    \item \textbf{Intelligent Routing Strategy}:
    \begin{itemize}
        \item Easy constraints → Direct solving (high confidence)
        \item Medium constraints → Conditional routing based on specific features
        \item Hard constraints → RL+LLM simplification (high confidence)
    \end{itemize}
\end{enumerate}

\textbf{Realistic Evaluation Methodology}

Unlike prediction-based approaches that require training data with known solving times, our complexity analyzer operates purely on constraint syntax, making it applicable in real-world scenarios where solving difficulty is unknown a priori.

\textbf{Results and Analysis}

Table~\ref{tab:complexity-routing} shows the distribution of constraints across complexity categories and the corresponding routing decisions.

\begin{table}[!t]
\centering
\caption{Complexity-guided routing system performance}
\label{tab:complexity-routing}
\begin{tabular}{lrrrrr}
\toprule
\textbf{Complexity} & \textbf{Count} & \textbf{\%} & \textbf{Avg. Score} & \textbf{Routing} & \textbf{Avg. Time (s)} \\
\midrule
Easy & 28,156 & 64.1\% & 2.8 & Direct & 12.4 \\
Medium & 12,447 & 28.3\% & 9.2 & Mixed & 45.7 \\
Hard & 3,311 & 7.5\% & 22.1 & RL+LLM & 127.3 \\
\bottomrule
\end{tabular}
\end{table}

The complexity-guided approach demonstrates several key advantages:

\begin{itemize}
    \item \textbf{Realistic Difficulty Identification}: 7.5\% of constraints are classified as hard based on syntactic complexity, correlating with significantly higher average solving times (127.3s vs 12.4s for easy constraints).
    
    \item \textbf{Effective Resource Allocation}: Hard constraints routed to RL+LLM show a 34\% improvement in success rate compared to direct solving, while easy constraints achieve optimal performance through direct solving.
    
    \item \textbf{Practical Applicability}: The routing system operates without requiring prior knowledge of solving difficulty, making it deployable in real-world symbolic execution scenarios.
\end{itemize}

\textbf{Comparative Performance Analysis}

Figure~\ref{fig:complexity-performance} compares the performance of our complexity-guided routing against uniform strategies (all direct solving vs. all RL+LLM).

\begin{itemize}
    \item \textbf{Overall Efficiency}: The hybrid approach achieves 23\% better overall performance than pure direct solving and 41\% better than applying RL+LLM to all constraints.
    
    \item \textbf{Resource Optimization}: By routing only 7.5\% of constraints to RL+LLM, the system minimizes computational overhead while maximizing benefits on truly challenging cases.
    
    \item \textbf{Scalability}: The static analysis approach scales linearly with constraint size, unlike prediction models that require expensive feature extraction and inference.
\end{itemize}

\begin{tcolorbox}[colback=white, colframe=black]
\textbf{Answer to RQ5:} Our complexity-guided routing system effectively identifies computationally challenging constraints using static analysis features, achieving practical deployment without requiring pre-solving knowledge. The system routes 7.5\% of constraints (those with high syntactic complexity) to RL+LLM simplification, resulting in 34\% improvement on hard constraints and 23\% overall system improvement compared to uniform direct solving. This demonstrates that our RL+LLM approach provides tangible benefits for challenging SMT constraints in realistic deployment scenarios.
\end{tcolorbox}
